\documentclass[11pt]{article}
\usepackage[margin=1in]{geometry}
\usepackage[T1]{fontenc}
\usepackage{lmodern}
\usepackage{microtype}
\usepackage{booktabs}
\usepackage{enumitem}
\usepackage{hyperref}
\usepackage{xurl}
\hypersetup{colorlinks=true,linkcolor=blue,citecolor=blue,urlcolor=blue}
\setlength{\parskip}{0.55em}
\setlength{\parindent}{0pt}

\title{ITM PC4 Melanocytic Signature and Frozen Route Markers\\
\large A Source-Identity and Overlap Audit}
\author{Inacio F. Vasquez\\Independent Researcher, United States}
\date{18 October 2026}

\begin{document}
\maketitle

\begin{abstract}
This note specifies an outcome-blind, reproducible comparison between the published in-transit melanoma (ITM) PC4 melanocytic transcriptional signature and the frozen seven-marker route-burden v2 feature definition. The current deliverable is an input-identity and signature-overlap audit, not a response score, prognostic model, or clinical analysis. The available evidence reconciles the reported cohort counts as 170 high-quality spatial regions of interest (ROIs) from 105 patients overall, including a described subset of 116 in-transit ROIs from 84 patients. The source matrix and complete PC4 signature were not retrieved for this audit; therefore no numerical overlap, correlation, survival estimate, or new prediction metric is claimed. These limitations are part of the result.
\end{abstract}

\section{Question and bounded result}

\textbf{Question.} Can the published PC4 melanocytic signature be compared with the frozen seven-marker route-burden feature definition without confusing cohort identity, signature identity, expression units, or clinical endpoints?

\textbf{Current result.}
\begin{quote}
\texttt{RESULT := REPRODUCIBLE INPUT AND OVERLAP GATE DEFINED; NUMERICAL OVERLAP NOT YET ESTABLISHED.}
\end{quote}

The audit is designed to verify matrix-to-manifest sample identity, count unique patients separately from ROIs, check signature completeness, inventory missing genes, and report exact and direction-labelled overlap between the two gene sets. It does not read expression values beyond the first-column gene inventory and sample header.

\section{Source and cohort accounting}

The published study reports 170 high-quality spatial ROIs from 105 distinct patients. It separately describes 116 ROIs from 84 patients as in-transit lesions, alongside ROIs from other metastatic sites. The latter count is a subset description and does not, by itself, demonstrate a contradictory dataset version. Any downstream patient-level analysis must identify the target population explicitly and join every ROI to its source-defined patient identifier.

The source record identifies a PC4 positive melanocytic signature of 36 genes and a PC4 negative signature of 56 genes. PMEL, TYRP1, TRPM1, GPR143, and TSPAN10 are identified among the high-contribution positive genes. PC4 is a PCA-derived transcriptional state; it is not interchangeable by definition with the frozen seven-marker route score.

Source records:
\begin{itemize}[leftmargin=1.5em]
  \item Article: \url{https://pmc.ncbi.nlm.nih.gov/articles/PMC13398476/}
  \item Dryad data record: \url{https://datadryad.org/dataset/doi:10.5061/dryad.kh18932pf}
  \item Related data/software record: \url{https://zenodo.org/records/19104096}
  \item GEO series: \url{https://www.ncbi.nlm.nih.gov/geo/query/acc.cgi?acc=GSE325123}
  \item Released analysis code: \url{https://github.com/ziyinhuang96/stgemox_itm_mel_tma/releases/tag/submission-v2}
\end{itemize}

The available ITM matrix is described as log2 counts per million (log2CPM). The frozen route-burden v2 empirical cumulative distribution function (ECDF) reference was built from a GSE78220 FPKM matrix. These units and preprocessing pipelines are not presumed interchangeable. The current audit intentionally does not apply the frozen ECDF reference to log2CPM data.

\section{Reproducible input contract}

The audit script \texttt{infra/ci/audit\_itm\_pc4\_route\_overlap.py} expects three local inputs:

\begin{enumerate}[leftmargin=2em]
  \item A gene-by-ROI CSV or TSV matrix. The first column contains unique gene symbols; remaining column headers are ROI/sample IDs. Expression values are not used.
  \item A complete PC4 signature table with columns \texttt{gene\_symbol,pc4\_direction}; direction must be \texttt{positive} or \texttt{negative}. The source-defined complete signature is expected to contain 36 positive and 56 negative genes.
  \item A sample manifest with columns \texttt{sample\_id,patient\_id}, one row per ROI and the source-defined patient identifier.
\end{enumerate}

Example invocation:

\begin{verbatim}
python3 infra/ci/audit_itm_pc4_route_overlap.py \
  --matrix ITM_log2CPM.csv \
  --signature pc4_signature.csv \
  --manifest itm_roi_manifest.csv \
  --output itm_pc4_route_overlap_receipt.json
\end{verbatim}

The JSON receipt records input paths, matrix/manifest identity, ROI and unique-patient counts, repeated ROIs per patient, signature dimensions, genes missing from the matrix, exact route-marker overlap, and direction-consistent versus direction-conflicting overlap. A size-gate override is appropriate only if a documented source correction establishes different signature dimensions; the source evidence should be recorded before changing expected counts.

\section{Comparison rule}

The frozen route-marker directions are:
\begin{center}
\begin{tabular}{ll}
\toprule
\textbf{Marker} & \textbf{Frozen route direction}\\
\midrule
KDM5B & high\\
CD36 & high\\
MITF & high\\
MTHFD2 & high\\
NGFR & high\\
NT5E & high\\
SOX10 & low\\
\bottomrule
\end{tabular}
\end{center}

For this audit only, a high-direction route marker is direction-consistent with a PC4-positive gene; the low-direction SOX10 marker is direction-consistent with a PC4-negative gene. This is a declared bookkeeping convention for comparison, not a claim that the PC4 and route signatures measure the same biological construct.

\section{Interpretation and failure rules}

\begin{itemize}[leftmargin=1.5em]
  \item Exact gene overlap is descriptive and does not establish biological equivalence.
  \item A matrix/manifest identity mismatch or incomplete PC4 signature blocks interpretation; unmatched rows must not be silently dropped.
  \item ROIs are not independent patients. Any survival analysis must use patient-level outcomes and account for repeated ROIs within patients.
  \item The frozen FPKM-based ECDF score must not be applied to the log2CPM matrix. A future same-pipeline reconstruction requires a separate, outcome-blind compatibility gate.
  \item A published survival association for PC4 is prognostic evidence for that signature; it does not validate the frozen seven-marker score as an anti-PD-1 response predictor.
\end{itemize}

\section{Treatment context: metadata gate, not a new score}

The audit now inventories optional source-defined manifest fields:
\texttt{treatment\_class}, \texttt{exposure\_status},
\texttt{biopsy\_timing}, \texttt{response\_endpoint}, and
\texttt{response\_status}. For each field, the JSON receipt records column
presence, reported and unresolved row counts, and counts by recorded category.
Blank or explicitly unknown values remain unresolved; treatment assignment,
timing, and outcomes are not inferred.

This is a treatment-context parameterization of the audit metadata, not a new
term in the frozen score and not a fitted clinical model. The seven-marker
feature definition, marker directions, preprocessing, and thresholds remain
unchanged. The audit does not estimate treatment-response associations or
score-by-treatment interactions.

A future interaction analysis requires source-supported patient-level linkage,
a defined pretreatment baseline, documented treatment class and exposure,
a prespecified endpoint and follow-up, and adequate treatment-comparison data.
Repeated ROIs must not be treated as independent patients. Even a detectable
score-by-treatment interaction would not alone establish treatment benefit or
support treatment selection; confounding and study design remain material.

\begin{quote}
\texttt{TREATMENT\_CONTEXT := SOURCE-RECORDED METADATA ONLY.}

\texttt{TREATMENT\_RESPONSE\_MODEL\_FIT := FALSE.}

\texttt{SCORE\_X\_TREATMENT\_INTERACTION\_TESTED := FALSE.}

\texttt{FROZEN\_SEVEN\_MARKER\_SCORE\_CHANGED := FALSE.}
\end{quote}

\texttt{BOUNDARY := $\neg$ established(treatment\_predictive\_validity)}

\section{Current evidence boundary}

The source matrix and complete PC4 signature were not successfully retrieved for the current audit. As a result, this report does not provide numerical overlap, correlation, patient-level survival estimates, or a new prediction metric. The script and regression tests define a reproducible gate for when the complete inputs are available; they do not substitute for those inputs.

\begin{quote}
\texttt{SOURCE\_COHORT\_ACCOUNTING := 170 ROI / 105 patients overall; 116 ROI / 84 patients in the described ITM subset.}

\texttt{PC4\_SIGNATURE\_EQUIVALENCE\_TO\_FROZEN\_ROUTE\_SCORE := NOT ESTABLISHED.}

\texttt{FROZEN\_V2\_SCORING\_ON\_LOG2CPM := NOT PERMITTED BY CURRENT COMPATIBILITY EVIDENCE.}

\texttt{NEW\_RESPONSE\_PREDICTION\_VALIDATION := NOT ESTABLISHED.}
\end{quote}

\texttt{BOUNDARY := $\neg$ established\_equivalence(PC4\_melanocytic\_signature, frozen\_7\_marker\_route\_score)}

\texttt{BOUNDARY := $\neg$ established\_response\_prediction(frozen\_7\_marker\_route\_score)}

\section{Final status}

This is a bounded source and methods note. It documents a reproducible comparison procedure and explicit stop conditions; it does not report a completed numerical comparison or expand the clinical claims.

\begin{thebibliography}{9}
\bibitem{ITMPC4}
Huang et al.
\newblock ITM melanoma spatial-transcriptomic study and PC4 melanocytic signature.
\newblock Article record: \url{https://pmc.ncbi.nlm.nih.gov/articles/PMC13398476/}.
\newblock Data and code records are linked in Section 2.
\end{thebibliography}



\\section{Source-code follow-up: existing ICI interaction specification}

Inspection of the released analysis script found an exploratory interaction section in \\url{https://github.com/ziyinhuang96/stgemox_itm_mel_tma/blob/submission-v2/geomx_wta_mel_tma_review.Rmd}, lines 4081--4152. The script reads a separate \\texttt{ITM\_treatment\_annotation.csv}, maps \\texttt{Unique.ID} to the patient key, joins an \\texttt{ICI} field, stratifies samples using the cohort PC4 survival cutpoint, and specifies a Cox model containing ICI, PC4 group, and their interaction for melanoma-specific survival. This establishes that the released code specifies an interaction analysis; it does not establish its numerical result, external validity, or causality.

The treatment annotation file and fitted output were not present in the inspected public code release. The expression matrix is listed on Dryad at \\url{https://datadryad.org/dataset/doi:10.5061/dryad.kh18932pf}, but the file endpoint returned HTTP 403 through the available retrieval route. The audit recognizes the source's literal \\texttt{ICI} column as an exposure-status alias and preserves recorded values without inferring regimen, timing, or response.

This is an observational survival interaction, not a randomized treatment-effect estimate. The cohort-derived PC4 cutpoint also prevents treating this code path as independent validation of a locked predictive threshold. Reanalysis requires the treatment annotation, patient linkage, endpoint data, complete expression matrix and signature, and a prespecified analysis plan.

\\begin{quote}
\\texttt{SOURCE\_CODE\_SPECIFIES\_ICI\_X\_PC4\_COX\_INTERACTION := TRUE.}\\par
\\texttt{INTERACTION\_NUMERICAL\_RESULT\_INDEPENDENTLY\_VERIFIED := FALSE.}\\par
\\texttt{CAUSAL\_TREATMENT\_BENEFIT\_ESTABLISHED := FALSE.}
\\end{quote}

\\texttt{BOUNDARY := $\\neg$ established(causal\_treatment\_benefit)}

\end{document}
